\documentclass[11pt]{article}
\usepackage[a4paper,margin=1in]{geometry}
\usepackage{fontspec}
\usepackage[boldfont=false]{xeCJK}
\setCJKmainfont{FandolSong-Regular.otf}
\usepackage{amsmath}
\usepackage{amssymb}
\usepackage{array}
\usepackage{booktabs}
\usepackage{graphicx}
\usepackage{adjustbox}
\usepackage{enumitem}
\setlist[itemize]{topsep=2pt,partopsep=0pt,parsep=0pt,itemsep=2pt,leftmargin=1.4em}
\setlist[enumerate]{topsep=2pt,partopsep=0pt,parsep=0pt,itemsep=2pt,leftmargin=1.6em}
\usepackage{parskip}
\usepackage{fvextra}
\fvset{breaklines=true,breakanywhere=true,fontsize=\small}
\setlength{\parindent}{0pt}

\begin{document}

\begin{center}
  {\LARGE\bfseries Nitrogen compounds}\\[0.35em]
  {\large A Level Chemistry}
\end{center}
\vspace{0.6em}

\section*{Primary and secondary amines}

An \textbf{amine} 胺 has an $\text{–NH}_2$ (primary) or $\text{–NH}$ (secondary) group.

\subsection*{Making amines}

\begin{itemize}
\item a \textbf{halogenoalkane} 卤代烷 with $\text{NH}_3$ in ethanol, heated under pressure → a primary amine.

\item a halogenoalkane with a primary amine, heated under pressure → a secondary amine.

\item \textbf{reduction} 还原 of an \textbf{amide} 酰胺 with $\text{LiAlH}_4$.

\item reduction of a \textbf{nitrile} 腈 with $\text{LiAlH}_4$ or $\text{H}_2/\text{Ni}$.

\end{itemize}

An amine (or \textbf{ammonia} 氨) reacts with an acyl chloride in a \textbf{condensation} 缩合 reaction to give an amide. The reagent is an \textbf{acyl chloride} 酰氯.

\subsection*{Basicity of amines}

The \textbf{basicity} 碱性 of an amine comes from the lone pair on its nitrogen, which can accept an $\text{H}^+$ from water.

\section*{Phenylamine and azo compounds}

\subsection*{Making phenylamine}

Make \textbf{phenylamine} 苯胺 from benzene in two steps: nitrate benzene to nitrobenzene, then reduce it with hot $\text{Sn}$ and concentrated $\text{HCl}$, followed by $\text{NaOH(aq)}$.

\subsection*{Reactions}

\begin{itemize}
\item with bromine water at room temperature → 2,4,6-tribromophenylamine (a white precipitate). The ring is activated, like phenol.

\item with $\text{HNO}_2$ (from $\text{NaNO}_2$ and dilute acid) below $10\,°\text{C}$ → a diazonium salt; warming this with water gives phenol.

\end{itemize}

\subsection*{Relative basicity}

\[
\text{phenylamine} < \text{ammonia} < \text{ethylamine}
\]

Ethylamine is the strongest base: its alkyl group pushes electron density onto the nitrogen, making the lone pair more available. Phenylamine is the weakest, because its lone pair is \textbf{delocalised} 离域 into the benzene ring, so it is less available to accept an $\text{H}^+$.

\subsection*{Azo compounds}

A \textbf{diazonium salt} 重氮盐 (benzenediazonium chloride) couples with \textbf{phenol} 苯酚 in $\text{NaOH(aq)}$ to form an \textbf{azo compound} 偶氮化合物. The azo group is $\text{–N=N–}$. Azo compounds are brightly coloured and are often used as \textbf{dye} 染料s; many other azo dyes are made the same way.

\section*{Amides}

An amide is made from ammonia or a primary amine with an acyl chloride at room temperature.

Its reactions:

\begin{itemize}
\item \textbf{hydrolysis} with aqueous acid or alkali, giving the carboxylic acid (or its salt) and the amine (or ammonium).

\item \textbf{reduction} of the C=O group with $\text{LiAlH}_4$ to give an amine.

\end{itemize}

An amide is a \textbf{much weaker base} than an amine, because the nitrogen lone pair is delocalised onto the neighbouring C=O group, so it is not available to accept an $\text{H}^+$.

\section*{Amino acids}

An \textbf{amino acid} 氨基酸 has both a basic $\text{–NH}_2$ group and an acidic $\text{–COOH}$ group.

\subsection*{Zwitterions and the isoelectric point}

The $\text{–COOH}$ can give its $\text{H}^+$ to the $\text{–NH}_2$ in the same molecule, forming a \textbf{zwitterion} 两性离子 ($\text{H}_3\text{N}^+\text{–CHR–COO}^-$) — an ion with both a positive and a negative end but no overall charge.

\begin{itemize}
\item in acid (low pH), the amino acid gains $\text{H}^+$ and becomes \textbf{positive}.

\item in alkali (high pH), it loses $\text{H}^+$ and becomes \textbf{negative}.

\item at one special pH, the \textbf{isoelectric point} 等电点, it is mostly the zwitterion with no net charge.

\end{itemize}

\subsection*{Peptide bonds}

Two amino acids join in a condensation reaction: the $\text{–COOH}$ of one and the $\text{–NH}_2$ of the other react, losing water and forming a \textbf{peptide bond} 肽键 (an amide link). Two amino acids give a \textbf{dipeptide} 二肽, three give a tripeptide.

\subsection*{Electrophoresis}

In \textbf{electrophoresis} 电泳, a mixture is placed in an electric field at a chosen pH:

\begin{itemize}
\item above its isoelectric point, an amino acid is negative and moves to the positive electrode.

\item below its isoelectric point, it is positive and moves to the negative electrode.

\item at its isoelectric point, it does not move. So different amino acids separate.

\end{itemize}


\end{document}
